\section{Background and Study Framing}

\subsection{Proof Maintenance Under Evolution}

In this paper, \emph{proof maintenance} means preserving or restoring successful verification across nearby program versions rather than proving each version from scratch. That framing matters because verified software evolves alongside the rest of a codebase: implementations change, helper abstractions move, and proof-relevant artifacts must co-evolve with the source program~\cite{Zaidman2010CoEvolutionProductionTest}. We distinguish three notions throughout the paper. \emph{Source evolution} refers to the program change itself. \emph{Proof brittleness} refers to the failure of previously successful verification under a nearby-version edit, even when the intended behavior remains unchanged. \emph{Semantic drift} refers instead to cases where the meaning of the program or its intended specification has genuinely changed. Keeping those notions separate matters empirically: without that distinction, a benchmark could blur maintenance failures together with genuine redesign. Prior work on proof refactoring shows that maintenance of proof artifacts is a legitimate research problem in its own right~\cite{Whiteside2011ProofScriptRefactoring,Whiteside2013RefactoringProofs}; this paper focuses on the empirical characterization of that problem in an auto-active verification setting rather than on automated repair.

\subsection{Why Dafny}

Dafny is an appropriate setting for this study because it makes proof-relevant structure visible at the source level while still relying on largely automatic proof checking~\cite{Leino2010Dafny}. Contracts, assertions, loop invariants, and ghost code appear directly in the program text, so maintenance effects are observable without requiring access to low-level interactive proof objects. At the same time, successful Dafny verification still depends on structured proof guidance and recurring verification idioms~\cite{Grov2016VerificationPatternsDafny}. Recent empirical work also suggests that usability and maintenance effort remain central to the practical cost of formal verification adoption in Dafny~\cite{Mugnier2025ImpactVerification}. This makes Dafny especially suitable for a benchmark whose purpose is not merely to count failures, but to examine which kinds of source-visible proof structure are most brittle under evolution.

\subsection{Semantics-Preserving Refactoring in This Paper}

The phrase \emph{semantics-preserving refactoring} is used narrowly in this paper. Following the general refactoring literature~\cite{Mens2004RefactoringSurvey,Kim2012RefactoringInTheWild,Tsantalis2014RefactoringOperations}, we study transformations intended to preserve program behavior while changing source structure. In this paper, that definition is further restricted to refactorings that may perturb verification outcomes by moving, reshaping, or redistributing proof-relevant structure in already verified Dafny programs. This excludes verifier bugs, deliberate semantic changes, intentionally weakened specifications, and unconstrained real-history edits. The point of controlled semantics-preserving refactoring here is not to model all software evolution, but to create nearby-version pairs with clean enough ground truth that proof brittleness can be measured reproducibly and compared across baselines.

\subsection{Study Objective and Evaluation Lens}

This paper studies how often verification breaks under controlled semantics-preserving refactoring, how those failures cluster by proof-relevant failure mode, and how much simple baseline recovery can restore. Dataset~A is therefore treated as a maintenance-centred benchmark rather than a plain corpus of transformed programs: each item is a nearby-version pair annotated with refactoring class, verification outcome, and baseline-recovery information. The evaluation lens is correspondingly simple and explicit: benchmark composition, breakage rate, baseline recovery, and the resulting residual repair gap. Repair and proof-automation work provide useful comparison points for baseline design~\cite{Gazzola2017AutomaticSoftwareRepair,First2020TacTok,SanchezStern2023Passport}, but the paper does not attempt to solve the repair problem itself. Instead, it defines the empirical problem sharply enough that Section~3 can introduce Dataset~A as a reusable substrate for studying nearby-version proof breakage under software evolution.
